\superchapter{Нэр томьёоны тайлбар}

{\small
\renewcommand{\arraystretch}{1.2}
\begin{longtable}{P{0.22} P{0.20} P{0.10} P{0.40}}
\toprule
\textbf{Англи нэр томьёо} & \textbf{Монгол нэршил} & \textbf{Товчлол} & \textbf{Тайлбар} \\
\midrule
\endfirsthead
\toprule
\textbf{Англи нэр томьёо} & \textbf{Монгол нэршил} & \textbf{Товчлол} & \textbf{Тайлбар} \\
\midrule
\endhead

% === Судалгааны суурь ойлголтууд ===
Computer Vision & Компьютерийн хараа & CV & Дижитал зураг, видеоноос утга бүхий мэдээлэл автоматаар гарган авах хиймэл оюуны дэд салбар. \\
Object Detection & Объект илрүүлэлт &  & Зураг доторх объектыг ангилал болон байршлын хамт тодорхойлох компьютерийн харааны даалгавар. \\
Retail Shelf Monitoring & Жижиглэн худалдааны лангууны мониторинг &  & Дэлгүүрийн лангуу дээрх барааны байршил, тоо, зөрчлийг технологийн аргаар хянах үйл ажиллагаа. \\
Planogram & Лангууны байршлын төлөвлөлт &  & Барааг лангуун дээр ямар дараалал, тоо, байрлалтай өрөхийг заасан стандарт төлөвлөгөө. \\
Fast-Moving Consumer Goods & Хурдан эргэлттэй өргөн хэрэглээний бараа & FMCG & Богино хугацаанд их хэмжээгээр эргэлтэд ордог өргөн хэрэглээний барааны ангилал. \\
\midrule

% === YOLO ба загварын архитектур ===
You Only Look Once & YOLO объект илрүүлэлтийн архитектур & YOLO & Объектыг нэг дамжуулалтаар ангилал ба байршлын хамт таамагладаг нэг шатлалт илрүүлэлтийн архитектурын гэр бүл. \\
YOLOv8 & YOLOv8 загвар & YOLOv8 & Энэ судалгаанд бараа илрүүлэх үндсэн хөдөлгүүр болгон ашигласан Ultralytics-ийн anchor-free объект илрүүлэгч загвар. \\
Anchor-Free Detector & Анкергүй илрүүлэгч &  & Урьдчилан тогтоосон anchor box ашиглахгүйгээр объектын байрлал, ангиллыг шууд таамагладаг илрүүлэлтийн архитектур. \\
Transfer Learning & Шилжүүлэн суралцах арга &  & Урьдчилан сурсан загварын мэдлэгийг шинэ, бага хэмжээний өгөгдөлтэй даалгаварт дахин ашиглах сургалтын арга. \\
Fine-Tuning & Нарийвчилсан дахин сургалт &  & Урьдчилан сурсан загварын параметрүүдийг тухайн домэйны өгөгдөл дээр дахин тохируулан шинэчлэх үйл явц. \\
Backbone & Суурь онцлог олборлогч сүлжээ &  & Зурагнаас анхан ба дунд түвшний шинж тэмдгийг гарган авах нейрон сүлжээний үндсэн хэсэг. \\
Neck & Онцлог нэгтгэх давхарга &  & Олон түвшний шинж тэмдгийг нэгтгэн detection head рүү дамжуулах архитектурын хэсэг. \\
Detection Head & Илрүүлэлтийн толгой хэсэг &  & Ангилал, bounding box, confidence зэрэг эцсийн таамгийг үүсгэдэг загварын гаралтын хэсэг. \\
Inference & Илрүүлэлтийн гүйцэтгэл &  & Сурсан загварыг ашиглан шинэ зураг дээр таамаглал гаргах ажиллагаа. \\
\midrule

% === Өгөгдөл ба сургалт ===
Dataset & Өгөгдлийн багц &  & Загвар сургах, баталгаажуулах, туршихад ашиглах дүрс болон шошгын цогц. \\
Annotation & Тэмдэглэгээ &  & Зураг дахь объектын ангилал, байршлыг шошголж өгсөн хүний бэлтгэсэн лавлагаа өгөгдөл. \\
Data Augmentation & Өгөгдөл нэмэгдүүлэлт &  & Эргүүлэх, гэрэл өөрчлөх, mosaic зэрэг хувиргалтаар сургалтын өгөгдлийн олон янз байдлыг нэмэх арга. \\
Training Set & Сургалтын багц &  & Загварын параметрийг шинэчлэхэд шууд ашигладаг өгөгдлийн дэд багц. \\
Validation Set & Баталгаажуулалтын багц &  & Сургалтын явцад загварын ерөнхийшүүлэх чадварыг хянахад ашигладаг өгөгдлийн дэд багц. \\
Test Set & Туршилтын багц &  & Эцсийн гүйцэтгэлийг бие даасан байдлаар үнэлэхэд ашигладаг өгөгдлийн дэд багц. \\
Epoch & Эпох &  & Сургалтын өгөгдлийн бүх жишээг нэг бүтэн удаа загвараар дамжуулсан мөчлөг. \\
Batch Size & Багцын хэмжээ &  & Нэг шинэчлэлтийн алхамд зэрэг боловсруулагдах сургалтын жишээний тоо. \\
Learning Rate & Суралцах хурд &  & Алдааны дагуу параметрийг ямар хэмжээгээр шинэчлэхийг заадаг гол гиперпараметр. \\
Loss Function & Алдааны функц &  & Загварын таамаг ба бодит шошгын зөрүүг тоон утгаар хэмжих математик функц. \\
Complete Intersection over Union Loss & Бүрэн давхцлын алдаа & CIoU loss & Bounding box-ийн давхцал, төвийн зай, харьцааны зөрүүг хамтад нь харгалзан оновчлох loss функц. \\
\midrule

% === Үнэлгээний хэмжүүрүүд ===
Bounding Box & Хязгаарлагч хүрээ &  & Илэрсэн объектыг зураг дээр тэгш өнцөгт хүрээгээр тэмдэглэсэн байршлын дүрслэл. \\
Intersection over Union & Давхцлын хувь & IoU & Таамагласан ба лавлагаа bounding box-ийн давхцлыг харьцуулан хэмждэг үзүүлэлт. \\
Mean Average Precision at IoU 0.50 & IoU=0.50 үеийн дундаж нарийвчлал & mAP@50 & Энэ судалгаанд объект илрүүлэгчийн гол үнэлгээний хэмжүүр болгон ашигласан дундаж нарийвчлалын үзүүлэлт. \\
Mean Average Precision from 0.50 to 0.95 & IoU 0.50--0.95 муж дахь дундаж нарийвчлал & mAP@50--95 & Илрүүлэлтийн чанарыг илүү хатуу босгуудаар нэгтгэн үнэлдэг нийлмэл хэмжүүр. \\
Precision & Нарийвчлалын хувь &  & Эерэг гэж таамагласан илрүүлэлтүүдээс үнэхээр зөв байсан хувийг илэрхийлдэг хэмжүүр. \\
Recall & Ололтын хувь &  & Бодит объектуудаас хэдийг нь зөв илрүүлснийг илэрхийлдэг хэмжүүр. \\
Confidence Score & Итгэлцлийн оноо &  & Загвар тухайн таамгийг зөв байх магадлал өндөр гэж үзэж буй түвшинг 0-1 хооронд илэрхийлсэн утга. \\
Confidence Threshold & Итгэлцлийн босго &  & Бага магадлалтай таамгийг шүүхийн тулд confidence score-д тавьдаг доод хязгаар. \\
False Positive & Хуурамч эерэг илрүүлэлт & FP & Байхгүй объектыг байгаа мэтээр буруу илрүүлсэн тохиолдол. \\
False Negative & Хуурамч сөрөг алдаа & FN & Байгаа объектыг илрүүлж чадаагүй тохиолдол. \\
Latency & Саатлын хугацаа &  & Хүсэлт илгээгдсэнээс хариу буцаж ирэх хүртэлх нийт хугацаа. \\
\midrule

% === Аудитын логик ба бизнесийн үнэлгээ ===
Audit Session & Аудитын сеанс &  & Нэг лангуу эсвэл нэг хяналтын цэг дээр зураг авч, илрүүлэлт хийж, үр дүн гаргасан бүртгэлт үйл явц. \\
Expected Count & Хүлээгдэж буй тоо хэмжээ &  & Планограм эсвэл барааны бүртгэлээр тухайн байршилд байх ёстой гэж тогтоосон тоо хэмжээ. \\
Detected Count & Илэрсэн тоо хэмжээ &  & Загварын inference-ийн дүнд тухайн ангиллын бараанаас бодитоор тоологдсон тоо. \\
Match Rate & Тохирлын хувь &  & Илэрсэн ба хүлээгдэж буй барааны нийцлийг хувиар илэрхийлсэн аудитын үндсэн үзүүлэлт. \\
Audit Status & Аудитын төлөв &  & Match rate болон зөрүүний дүрмээс хамааран PASS, WARNING, FAIL хэлбэрээр гарсан үнэлгээ. \\
Threshold & Шийдвэрийн босго &  & Аудитын төлөвийг ангилахдаа ашигладаг урьдчилан тогтоосон тоон хязгаар. \\
Stock-out & Лангуу дахь барааны тасалдал &  & Планограммаар байх ёстой бараа лангуун дээр байхгүй болсон нөхцөл. \\
Discrepancy & Тооллогын зөрүү &  & Хүлээгдэж буй болон илэрсэн тоо хэмжээний хооронд үүссэн ялгаа. \\
\midrule

% === Архитектур ба ажиллах орчин ===
Application Programming Interface & Хэрэглээний програмчлалын интерфейс & API & Клиент ба серверийн хооронд өгөгдөл солилцох дүрэм, төгсгөлийн цэгүүдийн багц. \\
Representational State Transfer API & REST төрлийн API & REST API & HTTP арга ба нөөцөд тулгуурлан зохион байгуулсан веб үйлчилгээний интерфейс. \\
JSON Web Token & JSON веб токен & JWT & Нэвтрэлт болон эрхийн мэдээллийг токен хэлбэрээр дамжуулж баталгаажуулах арга. \\
Cross-Origin Resource Sharing & Эх үүсвэр хоорондын хүсэлтийн бодлого & CORS & Өөр домэйнээс ирэх веб хүсэлтийг зөвшөөрөх эсэхийг удирддаг аюулгүй байдлын механизм. \\
FastAPI & FastAPI фреймворк & FastAPI & Python хэл дээр асинхрон веб үйлчилгээ хөгжүүлэхэд ашигласан серверийн фреймворк. \\
Pydantic Model & Pydantic загвар &  & Request, response өгөгдлийн бүтэц, төрлийг баталгаажуулахад ашигладаг Python схемийн загвар. \\
Asynchronous Processing & Асинхрон боловсруулалт &  & I/O хүлээлтийн үед бусад хүсэлтийг түгжихгүйгээр зэрэгцээ урсгал зохион байгуулах арга. \\
ASGI & Асинхрон серверийн гарц интерфейс & ASGI & Python веб фреймворк ба серверийн хооронд асинхрон харилцаа тогтоох стандарт интерфейс. \\
Uvicorn & Uvicorn сервер & Uvicorn & FastAPI зэрэг ASGI аппликейшнийг ажиллуулахад ашигласан хөнгөн жинтэй сервер. \\
MongoDB & MongoDB өгөгдлийн сан & MongoDB & JSON-тэй төстэй баримт бичгийн бүтэц ашигладаг NoSQL өгөгдлийн сан. \\
Motor Async Driver & Motor асинхрон драйвер & Motor & MongoDB-тэй FastAPI-ийн event loop орчинд асинхрон байдлаар холбогдох Python драйвер. \\
Singleton Pattern & Нэг хувьт загварчлал &  & Загварын объект санах ойд нэг удаа үүсээд дахин ашиглагдах зохион байгуулалтын хэв маяг. \\
Cold-Start Latency & Анхны ачаалалтын саатал &  & Загвар эсвэл үйлчилгээ шинээр ачаалагдахад эхний хүсэлт дээр нэмэлтээр үүсэх хүлээлтийн хугацаа. \\
Concurrency & Зэрэгцээ хүсэлтийн зохицуулалт &  & Олон хэрэглэгчийн хүсэлтийг нэг орчинд хооронд нь саад учруулахгүйгээр боловсруулах чадвар. \\
Retry Mechanism & Давтан оролдлогын механизм &  & Сүлжээ эсвэл үйлчилгээний түр доголдлын үед хүсэлтийг дахин илгээх логик. \\
Timeout & Хүлээлтийн дээд хугацаа &  & Хүсэлт хэт удааширсан үед алдаа гэж үзэн таслах дээд хугацааны босго. \\
Health Check & Эрүүл мэндийн шалгалт &  & Үйлчилгээ ажиллаж байгаа эсэхийг автоматаар шалгах хяналтын механизм. \\
Docker Container & Docker контейнер & Docker & Аппликейшн болон түүний хамаарлыг тусгаарласан орчинд тогтвортой ажиллуулах савлалтын нэгж. \\
\midrule

% === Шинжилгээ ба зохиомжийн арга хэрэгсэл ===
Use Case Diagram & Хэрэглээний тохиолдлын диаграмм &  & Оролцогч ба системийн үндсэн харилцан үйлчлэлийг дүрслэн үзүүлдэг UML диаграмм. \\
RACI Matrix & RACI хариуцлагын матриц & RACI & Үйл ажиллагаа бүр дээр хэн хариуцах, зөвшөөрөх, зөвлөлдөх, мэдээлэхийг ялган харуулдаг матриц. \\
Sequence Diagram & Дарааллын диаграмм &  & Объект, компонентуудын хооронд мессеж ямар дарааллаар дамжихыг хугацааны тэнхлэгээр харуулдаг UML диаграмм. \\
Component Diagram & Компонентын диаграмм &  & Системийн үндсэн бүрэлдэхүүн хэсэг, интерфейс, харилцан холбоог дүрслэн үзүүлдэг архитектурын диаграмм. \\
Deployment Diagram & Байршлын диаграмм &  & Системийн програм, сервер, өгөгдлийн сан, контейнерүүд ямар орчинд байрлаж ажиллахыг харуулдаг диаграмм. \\
MoSCoW Prioritization & MoSCoW эрэмбэлэлт & MoSCoW & Шаардлагыг Must, Should, Could, Won't ангиллаар тэргүүлэх ач холбогдлоор эрэмбэлэх арга. \\

\bottomrule
\end{longtable}
}
